Tôi, Kamishiro Ayaka, một tiểu thuyết gia 22 tuổi lừng danh với hàng
loạt những tác phẩm nổi tiếng là đi vào lòng người và được đánh giá rất
cao từ cả nhà sản xuất lẫn fan hâm mộ, giờ đây đang bị bí ý tưởng và
phải đi phượt một chuyến để mà có một cái gì đấy nảy ra trong đầu.

Đấy là suy nghĩ của tôi ban đầu, và tôi thật sự đã ngu ngốc mà tin rằng
đó là một phương án hay ho.

``Méo có...'' Tôi lẩm bẩm khi đang nằm dài ra trong phòng của mình ở nhà
trọ.

Méo có! Không có cái ý tưởng nào tràn vào đầu cả! Thậm chí bây giờ tôi
còn đang phải vật lộn với nhịp sinh hoạt khác hoàn toàn so với thường
ngày nữa, chả có thời gian đâu mà suy với chả nghĩ.

Ừ thì những địa điểm tôi ghé qua đều là mấy địa điểm khá nổi tiếng với
cảnh đẹp, lịch sử lâu đời các thứ, nhưng hoàn toàn chả có tác dụng gì
cả. Dù là một tiểu thuyết gia nhưng tôi thật sự chả thể cảm được nổi
khung cảnh hay lịch sử gì sất, thế nên tất nhiên là mấy thứ đấy chả thể
giúp được gì cho tôi rồi.

Nhưng đó chưa phải điểm tệ nhất...

Mà là, chuyến đi này tôi phải đi một mình! Nếu là trước đây thì không có
vấn đề gì. Nhưng hiện tại thì khác, tôi không thể sống nếu không có hơi
ấm quen thuộc của người yêu bên cạnh được!

``Serika ơi! Chị cô đơn quá!'' Tôi giãy đành đạch trong khi đang nằm,
khiến mền gối bị xáo trộn lung tung. ``Tại sao em vẫn còn chưa ra trường
chứ!''

Serika là một miêu nữ, hồi trước thì em ấy là thú cưng nhưng vào hai
tháng trước thì tôi có cầu hôn khi con bé khoảng 17 tuổi và hai đứa đã
chính thức thành người yêu với nhau. Đương nhiên thì chưa thành vợ chồng
được tại em ấy vẫn còn đang trong độ tuổi đến trường. À hay nên gọi mối
quan hệ của chúng tôi là vợ-vợ ấy nhỉ?

Mà thôi bỏ qua đi, tóm lại là cũng chính vì cái lý do đến trường đó mà
Serika không thể đồng hành cùng tôi trên chuyến hành trình tìm ra ý
tưởng cho cuốn tiểu thuyết tiếp theo được, nên là tôi mới phải đi một
mình đây.

Thậm chí tôi còn nhớ rõ câu nói của Serika khi mà nài nỉ em ấy đi theo
nữa cơ.

\emph{``Tuyệt đối không! Ngày mai em có bài học quan trọng!'' Serika chỉ
thẳng mặt tôi. ``Hơn nữa, nếu có em đi theo thì chắc chắn chị sẽ dành
toàn bộ thời gian để mà nựng em thôi chứ chả có tập trung vào tiểu
thuyết đâu!''}

Vậy đấy.

Ừ thì việc đó cũng không sai, nếu mà có Serika ở đây thì có khi tôi sẽ
dành toàn bộ thời gian của mình để mà cưng nựng con bé mất. Nhưng mà
việc không có em ấy cũng có giúp tôi tập trung thêm được tí nào đâu!

``Ahhhh...'' Tôi thở dài thườn thượt. ``Nhớ Serika quá đi...''

Bây giờ, chỉ cần được nghe giọng của con bé thôi cũng đủ mãn nguyện
rồi...

Bỗng nhiên cái điện thoại mà tôi vứt trên bàn phát ra tiếng chuông.

Ai gọi thế nhỉ?

``Serika!'' Tôi reo lên đầy vui mừng khi thấy cái tên được hiện lên trên
màn hình và ngay lập tức bắt máy. ``Serika đấy hả? Chị nhớ em quá!''

``Rồi rồi em biết, nghe cái giọng như cún con ấy nhỉ?'' Giọng của Serika
vang lên từ đầu dây bên kia có chút vang, là con bé đang tắm sao?

Oa, tưởng tượng ra cảnh đó thôi mà tâm trí của tôi còn phong phú hơn cả
mấy lúc nghĩ ý tưởng luôn mới ghê, đúng là không gì tuyệt vời bằng vợ
mình mà.

``Chị khỏe chứ?''

``Chị khỏe! Chị khỏe!'' Tôi ngồi bật dậy và hăng hái gật đầu. ``Chỉ là
có hơi thiếu hơi ấm của em thôi!''

``Là vẫn ổn đúng không?'' Serika nói bằng một giọng điệu ngán ngẩm mà
tôi đã quá quen thuộc ở em ấy. ``Chị đã có ý tưởng gì mới chưa?''

``Hehe, hỏi ngây thơ quá rồi đó Serika yêu dấu của chị.'' Tôi cười và
nói với một tông giọng nghe hết sức nguy hiểm, sau đó tuyên bố một cách
hùng hồn. ``Tất nhiên là chưa rồi!''

``...Em cúp máy đây.''

``A, đợi đã! Nói chuyện với chị thêm chút nữa đi mà!'' Tôi cố ra sức cầu
xin.

*

Đã khoảng hai ngày trôi qua kể từ khi chị Ayaka quyết định đi một chuyến
để mà tìm ý tưởng cho tiểu thuyết mới. Nhưng tôi nghĩ là nó cũng không
quá hiệu quả đâu khi mà chị ấy nếu không có tôi thì sẽ hoàn toàn vô dụng
và chả thể nghĩ một cái gì khác.

Ừ thì tôi biết là nếu không có tôi thì chị Ayaka sẽ mất hoàn toàn sức
sống rồi, nhưng mà kể cả tôi có mặt ở đó đi nữa thì tình hình chắc sẽ
không khả quan hơn là mấy đâu, mà có khi còn tệ hơn ấy chứ.

``Là mình nuông chiều chị ấy quá sao nhỉ?'' Tôi nằm dài trên ghế sofa và
suy nghĩ.

Đúng là tôi rất vui khi được chị Ayaka dựa vào. Và tôi cũng không phiền
để chị ấy cưng nựng hay là nấu cơm cho chị ấy. So với những gì mà chị
Ayaka làm cho tôi thì chừng này cũng chẳng đáng là bao. Nhưng nếu cứ thế
này thì tôi sợ chị Ayaka sẽ bị chê cười mất.

Thế nên để chị ấy luyện tập một chút cũng ổn, nó cũng có lợi một phần
cho tôi nữa.

Nhưng mà, tôi nghĩ bản thân cũng chả khá khẩm hơn là bao. Cứ mấy lúc nằm
một mình như thế này, trong đầu tôi lại tự tưởng tượng ra hình ảnh chị
Ayaka xoa đầu và đôi tai của mình. Cảm giác đó rất dễ chịu, và cũng rất
an tâm nữa.

Giờ lại không có chị ấy ở đây, tôi cũng thấy chút cô đơn.

Chỉ có một chút thôi...

``A! Mình nghĩ cái gì vậy!'' Tôi ngồi bật dậy, mặt đỏ bừng với mấy suy
nghĩ bật ra trong tình trạng yếu đuối vì thiếu hơi ấm của chị Ayaka.
``Thôi đi nấu ăn nào! Chị ấy cũng sắp về rồi!''

Tôi lật đật chạy thẳng vào bếp và bắt đầu lôi mấy thứ trong tủ lạnh ra.
Bật mí một chút là tài nấu nướng của tôi được chỉ dạy bởi chị Ayaka hồi
còn là một người sống độc lập. Nghe hơi khó tin nhỉ? Nhưng mà lúc đó chị
ấy thật sự làm việc nhà rất giỏi và nấu ăn cũng rất ngon. Đến cái mức mà
có thể chỉ dạy một đứa vụng về như tôi đây trở nên thành thạo chuyện bếp
núc như này cơ mà.

Tuy là bây giờ chị Ayaka khác hoàn toàn so với trước đây tại việc nhà
tôi làm hết rồi, nhưng nếu muốn thì chị ấy cũng có thể trở nên toàn năng
như trước thôi. Nhưng tôi hoàn toàn chẳng muốn chuyện đó xảy ra chút
nào. Tại nếu mà trở về như trước thì tôi làm sao có thể chăm sóc chị ấy
được nữa.

Hai năm trước, khi được chị Ayaka nhặt về, tôi đã rất ngưỡng mộ chị ấy.

Và tất nhiên, bây giờ vẫn vậy.

*

``Serika... chị về với em đây...'' Tôi bước về ngôi nhà thân yêu của
mình với dáng vẻ trông không khác gì một con zombie.

Chuyến đi của tôi kéo dài khoảng hai ngày, và trong hai ngày đó thì buộc
tôi phải tìm ra ý tưởng cho tiểu thuyết tiếp theo. Và tất nhiên là tôi
đã thành công, sau hai ngày đầy khó khăn.

Tuy nhiên, ý tưởng thì có rồi, nhưng tôi vẫn phải cố lếch xác về để gặp
người vợ miêu nữ đang chờ tôi ở nhà cái đã rồi làm gì thì làm, việc
thiếu hơi ấm của em ấy trong suốt hai ngày qua thật sự khiến tôi sắp
chết rồi đấy.

``Cuối cùng cũng đến nhà...'' Tôi đứng trước cửa nhà của mình và mở nó
ra.

Ở bên trong, Serika đã đứng chờ tôi sẵn với cái tạp dề trên người, đôi
tai mèo nhỏ nhắn trên đầu cũng với cái đuôi đang vẫy tít phía sau thể
hiện sự vui mừng, phản bội lại hoàn toàn khuôn mặt đang cố ra vẻ điềm
tĩnh của em ấy.

Tôi nghĩ... cái công tắc ngớ ngẩn của tôi lại được bật lên rồi.

``Chị về rồ-''

Không để Serika nói hết câu, tôi đã bế con bé lên theo kiểu công chúa và
tiến dần về phòng ngủ. Em ấy thật sự rất nhẹ nên việc này cũng không quá
khó khăn.

Ngược lại, tôi còn thấy thích ấy chứ.

``A, nè! Chị cần phải ăn chút gì đó cái đã!'' Serika cố gắng vùng vẫy và
ngăn tôi lại, nhưng hiện giờ thì tôi chẳng còn nghe lọt một câu nào nữa
cả. ``Ít nhất thì phải đi tắm cái đã chứ!''

Mặc cho Serika đang vùng vẫy trên tay, tôi vẫn từ từ tiến gần đến phòng
ngủ.

*

Sau nụ hôn với Serika đang bị tôi đè ra giường trước mặt thì tôi cũng
lấy lại một chút tỉnh táo. Thế nhưng với cái gương mặt đỏ bừng, khóe mắt
rơm rớm nước và quần áo bị tôi gần như cởi hết ra thì cái lý trí cuối
cùng này của tôi có khi sẽ bay màu luôn mất.

Có lẽ, tôi nên làm nguội cái đầu một chút.

``Chị xin lỗi, quả nhiên chị nên đi tắm trước.'' Tôi nói và định ngồi
dậy.

``Không!''

``Ặc!'' Tôi kêu lên một tiếng trước cái ôm bất ngờ của Serika.

Đây là Serika ở dere mod đấy à? Dạo này hiếm khi thấy nó xuất hiện lắm.
Có lẽ không phải chỉ có mình tôi phải chịu đựng sự cô đơn trong suốt hai
ngày qua nhỉ?

``Con mèo tsundere này...'' Tôi khẽ vuốt ve cái đuôi vốn là điểm yếu của
Serika khiến con bé giật mình. ``Nếu cứ như vậy thì chị biết kiềm chế
kiểu gì đây...''

``Không cần phải kiềm chế đâu...'' Serika nằm bên dưới nhìn tôi với một
ánh mắt ngượng ngùng. ``Em cũng nhớ hơi ấm của chị...''

``Này là tự em nói đấy nhé.''

Vậy thì, tôi không khách sáo nữa vậy.

Một hồi cái tính tsun trở lại, Serika sẽ dễ thương lắm đây.
